\documentclass[11pt]{article}
\usepackage[margin=1in]{geometry}
\usepackage{amsmath,amssymb,amsfonts,mathtools}
\usepackage{array,booktabs,enumitem,graphicx,xcolor,hyperref}
\usepackage[T1]{fontenc}
\usepackage[utf8]{inputenc}
\title{SU(3) x SU(2) x U(1): The residual symmetry of extended conformal
  gravity}
\author{Source-backed arXiv sample}
\date{}
\begin{document}
\maketitle
\section{Extracted Lists}

The list environments below are extracted from arXiv LaTeX source and wrapped for pdfLaTeX validation.

\begin{enumerate}
         \item Conformal symmetry must underlie our mathematical
                    description of nature.
         \item Naive conformal gravity fails.
     \end{enumerate}

\begin{enumerate}
              \item Scale changes, generated by the dilation operator, $D$
directly produce the allowed choices of units by rescaling lengths locally.
              \item Conformal translations.  The four generators, $K_{a}$,
translate infinitesimal vectors in inverse coordinates.  Such a vector is
inverted through a unit sphere, translated, then re-inverted.  The
transformation clearly involves an arbitrary displacement vector.  The thing to
notice -- or rather not to notice -- is that the vector is also rotated and
rescaled.  Since both rotating and stretching are already independent
transformations, the group closes.
          \end{enumerate}
\end{document}
